%============================================================
% TEMPLATE PROPOSAL TESIS
% Program Studi Ilmu Komputer (S2)
% Fakultas Teknologi Informasi
% Universitas Nusa Mandiri
% Revisi 2 - Februari 2024
%============================================================

\documentclass[12pt, a4paper]{report}

%------------------------------------------------------------
% PACKAGES
%------------------------------------------------------------
\usepackage[utf8]{inputenc}
\usepackage[T1]{fontenc}
\usepackage{times}                  % Times New Roman
\usepackage[bahasa]{babel}          % Bahasa Indonesia
\usepackage{geometry}
\usepackage{setspace}
\usepackage{titlesec}
\usepackage{titletoc}
\usepackage{graphicx}
\usepackage{hyperref}
\usepackage{indentfirst}
\usepackage{enumitem}
\usepackage{tabularx}
\usepackage{booktabs}
\usepackage{fancyhdr}
\usepackage{caption}
\usepackage{cite}                   % IEEE citation style
\usepackage{notoccite}
\usepackage{amsmath}
\usepackage{microtype}

%------------------------------------------------------------
% MARGIN: Kiri-Atas-Kanan-Bawah = 4-3-2.5-2.5 cm
%------------------------------------------------------------
\geometry{
    left   = 4cm,
    top    = 3cm,
    right  = 2.5cm,
    bottom = 2.5cm,
    a4paper
}

%------------------------------------------------------------
% SPASI 1.5
%------------------------------------------------------------
\setstretch{1.5}

%------------------------------------------------------------
% FONT: Times New Roman 12pt (default via \usepackage{times})
%------------------------------------------------------------

%------------------------------------------------------------
% JUDUL BAB: Kapital, Tengah, Tebal, Tanpa Garis Bawah
%------------------------------------------------------------
\titleformat{\chapter}[block]
    {\centering\bfseries\large}
    {BAB \Roman{chapter}}
    {1em}
    {\MakeUppercase}
\titlespacing*{\chapter}{0pt}{0pt}{20pt}

% SUBBAB LEVEL 1 (1.1)
\titleformat{\section}
    {\bfseries\normalsize}
    {\thesection.}
    {0.5em}
    {}
\titlespacing*{\section}{0pt}{12pt}{6pt}

% SUBBAB LEVEL 2 (1.1.1)
\titleformat{\subsection}
    {\bfseries\normalsize}
    {\thesubsection.}
    {0.5em}
    {}
\titlespacing*{\subsection}{0pt}{10pt}{4pt}

% SUBBAB LEVEL 3 (1.1.1.1)
\titleformat{\subsubsection}
    {\normalfont\normalsize}
    {\thesubsubsection.}
    {0.5em}
    {}
\titlespacing*{\subsubsection}{0pt}{8pt}{2pt}

% Nonaktifkan penomoran bab di ToC (pakai huruf Romawi manual)
\renewcommand{\thechapter}{\Roman{chapter}}
\renewcommand{\thesection}{\arabic{chapter}.\arabic{section}}
\renewcommand{\thesubsection}{\arabic{chapter}.\arabic{section}.\arabic{subsection}}
\renewcommand{\thesubsubsection}{\arabic{chapter}.\arabic{section}.\arabic{subsection}.\arabic{subsubsection}}

%------------------------------------------------------------
% DAFTAR ISI
%------------------------------------------------------------
\addto\captionsbahasa{%
    \renewcommand{\contentsname}{DAFTAR ISI}%
    \renewcommand{\listfigurename}{DAFTAR GAMBAR}%
    \renewcommand{\listtablename}{DAFTAR TABEL}%
    \renewcommand{\bibname}{DAFTAR REFERENSI}%
    \renewcommand{\chaptername}{BAB}%
}

%------------------------------------------------------------
% NOMOR HALAMAN
%------------------------------------------------------------
\pagestyle{fancy}
\fancyhf{}
\fancyhead[R]{\thepage}
\renewcommand{\headrulewidth}{0pt}
\fancypagestyle{plain}{%
    \fancyhf{}
    \fancyfoot[C]{\thepage}
    \renewcommand{\headrulewidth}{0pt}
}

%------------------------------------------------------------
% HYPERREF (tautan PDF)
%------------------------------------------------------------
\hypersetup{
    colorlinks = true,
    linkcolor  = black,
    citecolor  = black,
    urlcolor   = black
}

%------------------------------------------------------------
% CAPTION TABEL & GAMBAR
%------------------------------------------------------------
\captionsetup[table]{position=above, labelsep=period, singlelinecheck=false}
\captionsetup[figure]{position=below, labelsep=period}

%============================================================
% DATA PROPOSAL — UBAH SESUAI KEBUTUHAN
%============================================================
\newcommand{\judulProposal}{OPTIMASI DINAMIS PEMODELAN NETWORK MARKOWITZ UNTUK MANAJEMEN PORTOFOLIO MATA UANG KRIPTO}
\newcommand{\namaMahasiswa}{Nama Lengkap Mahasiswa}
\newcommand{\nimMahasiswa}{XXXXXXXXX}
\newcommand{\tahunProposal}{2024}

%============================================================
% DOKUMEN MULAI
%============================================================
\begin{document}

%------------------------------------------------------------
% HALAMAN SAMPUL
%------------------------------------------------------------
\begin{titlepage}
    \centering
    \setstretch{1.5}

    {\bfseries\large \judulProposal \par}

    \vspace{2cm}
    \includegraphics[width=5cm]{logo_unm.png} % Ganti dengan path logo UNM

    \vspace{1.5cm}
    {\bfseries\Large PROPOSAL TESIS \par}

    \vspace{0.5cm}
    {\normalsize Diajukan sebagai salah satu syarat untuk memperoleh gelar \par}
    {\normalsize Magister Komputer (M.Kom) \par}

    \vspace{1.5cm}
    {\bfseries\large \namaMahasiswa \par}
    {\bfseries\large \nimMahasiswa \par}

    \vfill

    {\bfseries\normalsize Program Studi Ilmu Komputer (S2) \par}
    {\bfseries\normalsize Fakultas Teknologi Informasi \par}
    {\bfseries\normalsize Universitas Nusa Mandiri \par}
    {\bfseries\normalsize Jakarta \par}
    {\bfseries\normalsize \tahunProposal \par}
\end{titlepage}

%------------------------------------------------------------
% HALAMAN PENGESAHAN
%------------------------------------------------------------
\pagenumbering{roman}
\setcounter{page}{2}

\chapter*{HALAMAN PENGESAHAN}
\addcontentsline{toc}{chapter}{HALAMAN PENGESAHAN}

\begin{center}
    {\bfseries PROPOSAL TESIS}
\end{center}

\vspace{0.5cm}
\noindent
Proposal tesis ini diajukan oleh:

\vspace{0.5cm}
\begin{tabular}{lll}
    Nama          & : & \namaMahasiswa \\
    NIM           & : & \nimMahasiswa \\
    Program Studi & : & Ilmu Komputer \\
    Fakultas      & : & Teknologi Informasi \\
    Jenjang       & : & Strata Dua (S2) \\
    Judul Tesis   & : & \judulProposal \\
\end{tabular}

\vspace{2cm}
\noindent
Jakarta, \today

\vspace{1cm}
\begin{tabular}{cc}
    Dosen Pembimbing & \\[3cm]
    \underline{\hspace{8cm}} & \\
    Nama Dosen Pembimbing & \\
    NIDN. XXXXXXXXXX & \\
\end{tabular}

%------------------------------------------------------------
% DAFTAR ISI
%------------------------------------------------------------
\tableofcontents
\clearpage

%============================================================
% BAB I — PENDAHULUAN
%============================================================
\pagenumbering{arabic}
\setcounter{page}{1}

\chapter{PENDAHULUAN}

\section{Latar Belakang}

Mata uang kripto (\textit{cryptocurrency}) telah berkembang menjadi salah satu aset investasi digital yang sangat diminati namun memiliki tingkat volatilitas ekstrem. Dalam manajemen portofolio tradisional, model \textit{Mean-Variance} dari Markowitz kerap digunakan untuk mengalokasikan aset demi mencapai kombinasi \textit{return} dan risiko yang optimal. Sayangnya, model klasik ini sangat rentan terhadap *noise* dan estimasi matriks korelasi yang tidak stabil, terutama pada saat gejolak pasar (*market crash*) seperti fenomena \textit{crypto winter} \cite{ref1}.

Ketika terjadi guncangan pasar, sebagian besar aset kripto cenderung jatuh secara bersamaan, merusak struktur korelasi normal dan menyebabkan portofolio standar mengalami kerugian parah (*drawdown* yang dalam). Untuk mengatasi tantangan tersebut, penelitian terkini mulai menggabungkan Teori Jaringan Kompleks (\textit{Complex Network Theory}) ke dalam optimasi portofolio \cite{ref2}. Penggunaan instrumen seperti \textit{Minimum Spanning Tree} (MST) dan \textit{Eigenvector Centrality} terbukti efisien dalam memetakan interaksi antar aset dan menghukum (*penalize*) aset-aset yang menjadi titik pusat kegagalan sistemik. 

Kendati model \textit{Network Markowitz} statis menunjukkan proteksi yang lebih relevan dibandingkan \textit{Classical Markowitz}, penentuan faktor penalti sentralitas g ($\gamma$) yang kaku kerap menimbulkan masalah di fase pasar yang dinamis (misalnya fase \textit{recovery} atau \textit{bullish}). Oleh karena itu, diperlukan pendekatan adaptif yang mengintegrasikan teknik pembersihan sinyal berbasis \textit{Random Matrix Theory} (RMT) disandingkan dengan optimalisasi \textit{Grid Search} dinamis (\textit{rolling window}) agar portofolio dapat membentengi aset di saat *crypto winter* tanpa mengorbankan rasio *upside gain* di saat pembalikan arah (*bullish*).

\section{Identifikasi Masalah}

Berdasarkan latar belakang di atas, dapat diidentifikasi masalah sebagai berikut:
\begin{enumerate}[label=\alph*.]
    \item Model \textit{Classical Markowitz} rentan terhadap *spurious correlations* pada aset kripto yang bervolatilitas sangat tinggi, terutama pada kondisi krisis ekstrem.
    \item Implementasi \textit{Network Markowitz} yang telah ada sering kali menggunakan hiper-parameter penalti sentralitas ($\gamma$) bernilai konstan (statis), yang berpotensi menjadi bumerang saat pasar memasuki rezim *recovery* atau *bullish*.
    \item Kurangnya model sistematis yang secara metodis menyesuaikan struktur jaringan portofolio dengan rezim siklus guncangan harga terkini menggunakan teknik optimasi *rolling window* untuk aset-aset kripto utama.
\end{enumerate}

\section{Tujuan Penelitian}

Tujuan dari penelitian ini adalah:
\begin{enumerate}[label=\alph*.]
    \item Menganalisis keandalan metodologi \textit{Network Markowitz} (dengan integrasi RMT filter dan Eigenvector Centrality) dalam menekan ekstrimitas *downside risk* dibandingkan pendekatan portofolio naif dan konvensional.
    \item Merancang dan menguji model \textit{Network Markowitz} adaptif (\textit{Grid Search Optimization}) yang mampu melakukan re-kalibrasi dinamis dengan menggunakan paradigma \textit{rolling window} pada berbagai lanskap pasar (*Bearish*, *Recovery*, *Stable*).
    \item Membandingkan performa perlindungan risiko sistemik (VaR) dan asimetri imbal hasil (*Rachev Ratio*) antara pemodelan baru dengan metode-metode *baseline* pada reksadana aset kripto.
\end{enumerate}

\section{Ruang Lingkup Penelitian}

Ruang lingkup penelitian ini dibatasi pada:
\begin{enumerate}[label=\alph*.]
    \item Objek penelitian terfokus pada data fluktuasi harga harian dari 10 (\textit{sepuluh}) aset kripto utama dalam kerangka waktu historis termasuk masa resesi \textit{crypto winter} (14 September 2017 hingga 17 Oktober 2019).
    \item Metode yang dibandingkan secara teknis mencakup \textit{Equally Weighted} (EW), \textit{Classical Markowitz} (CM), \textit{Glasso Markowitz} (GM), \textit{Network Markowitz} statis ($\gamma=0, 1.0, 2.0$), serta *Optimized* \textit{Network Markowitz} secara dinamis berbasis \textit{Grid Search}.
    \item Pengujian (\textit{backtesting}) dilakukan dalam \textit{out-of-sample rolling window} (120 observasi ke belakang dengan frekuensi penyesuaian *rebablance* 7 hari) yang disimulasikan menggunakan \textit{transaction cost} atau estimasi biaya bursa (0.1\%).
\end{enumerate}

\section{Sistematika Penulisan}

Sistematika penulisan proposal tesis ini disusun sebagai berikut:

\textbf{BAB I PENDAHULUAN} \\
Bab ini membahas latar belakang penelitian, identifikasi masalah, tujuan penelitian, ruang lingkup penelitian, dan sistematika penulisan.

\textbf{BAB II LANDASAN/KERANGKA PEMIKIRAN} \\
Bab ini membahas kerangka teori yang relevan mencakup Teori Portofolio Modern, \textit{Complex Network Theory}, dan Manajemen Risiko, serta tinjauan pustaka terhadap penelitian terdahulu di bidang Robo-Advisory kripto.

\textbf{BAB III METODOLOGI PENELITIAN} \\
Bab ini menjelaskan alur sistematis eksperimen yang digunakan untuk memproses data instrumen kripto, penyaringan RMT pada matriks korelasi historis, pembangunan MST, fungsi objektif Markowitz modifikasi, dan komputasi skema \textit{backtesting}.

%============================================================
% BAB II — LANDASAN / KERANGKA PEMIKIRAN
%============================================================
\chapter{LANDASAN/KERANGKA PEMIKIRAN}

\section{Kerangka Teori}

\subsection{Modern Portfolio Theory (Markowitz)}
Teori Portofolio Modern, yang dipelopori oleh Harry Markowitz, berupaya memaksimalkan imbal hasil yang diharapkan (*expected return*) pada tingkat risiko (\textit{variance}) tertentu, atau sebaliknya. Masalah mendarnya adalah bahwa kovarians dari kumpulan asst finansial sangat sensitif terhadap nilai-nilai ekstrem historis, yang dikenal sebagai *Markowitz Curse* \cite{ref4}.

\subsection{Random Matrix Theory (RMT) dan Kompleksitas Jaringan}
Teori \textit{Random Matrix} memungkinkan disaringnya *noise* dari struktur korelasi dengan membandingkan nilai eigen (\textit{eigenvalues}) struktur empiris terhadap nilai batas distribusi teoritis Marchenko-Pastur \cite{ref3}. Ini merupakan vital element sebelum dilakukan visualisasi graf berupa \textit{Minimum Spanning Tree} (MST). 

\subsection{Network Markowitz}
Diperkenalkan baru-baru ini untuk penanganan robo-advisory pada kripto, komponen sentralitas *eigenvector* ditambahkan sebagai instrumen pinalti di dalam penyelesaian optimasi *Mean-Variance* \cite{ref1}. Sentralitas mewakili kerentanan sebuah aset mentransmisikan *shock* pada seluruh jaringan koin di pasar.

\subsection{Penelitian Terdahulu}

Beberapa penelitian terdahulu yang relevan dengan penelitian ini antara lain:

\begin{table}[h!]
    \centering
    \caption{Perbandingan Penelitian Terdahulu}
    \label{tab:penelitian_terdahulu}
    \begin{tabularx}{\textwidth}{|c|X|X|X|}
        \hline
        \textbf{No} & \textbf{Penulis / Tahun} & \textbf{Judul} & \textbf{Hasil} \\
        \hline
        1 & Giudici, et al. (2020) & Network Models to Improve Automated Cryptocurrency \dots & Mengusulkan Network Markowitz dan sukses mendemonstrasikan perbaikan struktur dibandingkan Markowitz biasa di era crypto winter. \\
        \hline
        2 & Papenbrock (2011) & Financial networks and risk management & Bukti bahwa penggunaan analisis jaringan meningkatkan akurasi *risk metrics*. \\
        \hline
        3 & Mantegna (1999) & Hierarchical structure in financial markets & MST sukses memotret topologi kedekatan dan konektivitas sektor finansial dari pergerakan deret waktu pasar saham. \\
        \hline
    \end{tabularx}
\end{table}

%============================================================
% BAB III — METODOLOGI PENELITIAN
%============================================================
\chapter{METODOLOGI PENELITIAN}

\section{Tahapan Penelitian}

Penelitian ini dilaksanakan melalui tahapan-tahapan sebagai berikut:

\begin{enumerate}
    \item \textbf{Studi Literatur} \\
    Melakukan kajian komprehensif terhadap jurnal yang membahas \textit{crypto portfolio optimization}, \textit{graph theory}, dan \textit{Network Markowitz}.

    \item \textbf{Pengumpulan Data \& Pengkondisian} \\
    Mengumpulkan seri harga \textit{close} \textit{cryptocurrency}. Harga lalu diubah menjadi format \textit{log returns}.

    \item \textbf{Pra-pemrosesan Data (RMT Filtering)} \\
    Data *returns* akan diubah menjadi Matriks Korelasi (Pearson). Selanjutnya struktur *noise* historis dipotong (di-nol/gantikan) menggunakan mekanisme filtrasi nilai *eigen* (*eigenvalue clipping boundary* batas Marchenko-Pastur).

    \item \textbf{Kuantifikasi Jaringan Aset (MST \& \textit{Centrality})} \\
    Pembangunan jarak konektivitas (*distance matrix*) dari hasil korelasi terfilter untuk diekstraksi ke bentuk graf pohon *Minimum Spanning Tree* (MST). *Node importance* kemudian diukur melalui \textit{Eigenvector Centrality}.

    \item \textbf{Optimasi Jaringan Adaptif (Dynamic Grid-Search)} \\
    Merancang algoritma komputasi untuk menyesuaikan parameter *impact* korelasi ($\gamma$) yang secara otomatis mengekang instrumen-instrumen bervolatilitas sistemik pada jendela uji terkalibrasi ke belakang (\textit{backtrack validity}).

    \item \textbf{Eksekusi \textit{Backtesting} Portofolio} \\
    Mensimulasikan pembelian pada titik waktu (*t*) dan meninjau portofolio secara periodik menggunakan sistem \textit{Rolling Window}. Terdapat perlakuan pengenaan \textit{slippage/transaction log} pada *rebalancing* harian.

    \item \textbf{Pengukuran Evaluasi Resiko (\textit{Performance Metrics})} \\
    Mengukur nilai profit kumulatif, asimetri VaR 95\%, hingga penyesuaian \textit{Sharpe Ratio} dan \textit{Rachev Ratio} di sepanjang berbagai transisi fasa pasar ekstrem.
\end{enumerate}

\section{Alat dan Bahan Penelitian}

\subsection{Perangkat Lunak}
\begin{itemize}
    \item Sistem Operasi: Windows 10/11
    \item Bahasa Pemrograman: Python 3.x (dengan ekosistem Anaconda)
    \item Framework/Library: Pandas, Numpy, Scipy (Optimization), Scikit-Learn, NetworkX (Graph Analytics).
\end{itemize}

\section{Dataset}

Dataset yang akan digunakan adalah sekumpulan (\textit{pool}) 10 aset berbasis *cryptocurrency* berkapitalisasi tinggi yang tercatat aktif diperdagangkan secara bersinambung. Terdiri dari *Bitcoin* (BTC), *Ethereum* (ETH), *Ripple* (XRP), *Litecoin* (LTC), dan lain sebagainya. Observasi direntangkan secara sengaja mencakup era penggelembungan (*speculative bubble*), era krisis berkepanjangan (\textit{crypto winter bear market}), hingga skema normalisasi *stable*.

\section{Metode/Algoritma yang Digunakan}

Optimasi yang diajukan akan merubah fungsi pencarian model klasik Markowitz menjadi kerangka berbasis *penalty function* dinamis. Formula \textit{Network Markowitz}:

\begin{equation}
    \min_{w} \left( w^T \cdot S_f \cdot w + \gamma \sum (C_e \cdot w) \right)
    \label{eq:networkmarkowitz}
\end{equation}

\noindent
Keterangan: \\
$w$ = Vektor alokasi bobot untuk setiap aset kripto (sumbangan = 1). \\
$S_f$ = Matriks Kovarians terfilter RMT. \\
$\gamma$ = Parameter skalar reguler untuk tingkat penghukuman (*penalty level*) sentralitas graf. \\
$C_e$ = Vektor skor *Eigenvector Centrality* tiap-tiap entitas *node*.

Pada penelitian ini, bobot $\gamma$ tidak akan dilakukan \textit{hard-coded statis}, melainkan secara luwes dan \textit{rolling} akan difungsikan optimasi *grid validation* berbasis metrik obyektif \textit{Sharpe Ratio} periode belakang.

\section{Rencana Jadwal Penelitian}

\begin{table}[h!]
    \centering
    \caption{Rencana Jadwal Penelitian}
    \label{tab:jadwal}
    \begin{tabularx}{\textwidth}{|c|X|c|c|c|c|c|c|}
        \hline
        \textbf{No} & \textbf{Kegiatan}
            & \textbf{Bln 1} & \textbf{Bln 2} & \textbf{Bln 3}
            & \textbf{Bln 4} & \textbf{Bln 5} & \textbf{Bln 6} \\
        \hline
        1 & Studi Literatur            & \checkmark & & & & & \\
        \hline
        2 & Pengumpulan Data           & \checkmark & \checkmark & & & & \\
        \hline
        3 & Pra-pemrosesan Data        & & \checkmark & \checkmark & & & \\
        \hline
        4 & Perancangan Model/Sistem   & & & \checkmark & \checkmark & & \\
        \hline
        5 & Implementasi               & & & & \checkmark & \checkmark & \\
        \hline
        6 & Pengujian dan Evaluasi     & & & & & \checkmark & \checkmark \\
        \hline
        7 & Penulisan Laporan          & & & \checkmark & \checkmark & \checkmark & \checkmark \\
        \hline
    \end{tabularx}
\end{table}

%============================================================
% DAFTAR REFERENSI (IEEE Style)
%============================================================
\renewcommand{\bibname}{DAFTAR REFERENSI}

\begin{thebibliography}{99}

\bibitem{ref1}
P. Giudici, A. Sariev, and G. Toscani, ``Network Models to Improve Automated Cryptocurrency Portfolio Management,'' \emph{Risks},
vol. 8, no. 3, p. 96, 2020. https://doi.org/10.3390/risks8030096.

\bibitem{ref2}
S. S. Momeni dan R. E. Rostami, ``Portfolio selection using a hybrid methodology of network theory and optimization,'' \emph{Annals of Operations Research},
vol. 308, pp. 317--345, 2021.

\bibitem{ref3}
V. A. Marchenko and L. A. Pastur, ``Distribution of eigenvalues for some sets of random matrices,'' \emph{Mathematics of the USSR-Sbornik},
vol. 1, no. 4, pp. 457--483, 1967.

\bibitem{ref4}
H. Markowitz, ``Portfolio Selection,'' \emph{The Journal of Finance},
vol. 7, no. 1, pp. 77--91, 1952.

\bibitem{ref5}
R. N. Mantegna, ``Hierarchical structure in financial markets,'' \emph{The European Physical Journal B},
vol. 11, no. 1, pp. 193--197, 1999.

\end{thebibliography}

\end{document}
